\documentclass[a4paper,12pt]{article}
\input{LaTeX-ToolKit/template.tex}
\begin{document}
\title{Band Theory}
\author{沈威宇}
\date{\temtoday}
\titletocdoc
\sct{Band Theory (能帶理論) and Semiconductor}
\ssc{Basics}
\sssc{Electronic band structure or band structure}
The electronic band structure (or simply band structure) of a solid describes the range of energy levels that electrons may or may not have within it.
\sssc{(Energy) band (能帶)}
Atomic orbitals from multiple atoms combine to form molecular orbitals, such as bonding and anti-bonding orbitals. When a vast number of atoms—typically on the order of $10^{20}$ or more—come together to form a solid, the energy differences between these molecular orbitals become extremely small. Consequently, a cluster of closely spaced energy levels can be treated as a continuous energy band.
\sssc{Fermi level (費米能階)}
The Fermi level of a solid-state body is the thermodynamic work required to add one electron to the body. For Cu, $E_F=7.0$ eV.
\sssc{Work function (功函數)}
For a material with Fermi level $E_F$ and energy of an electron infinitely far in vacuum $E_{\tx{vac}}$, the work function is
\[W=E_{\text{vac}}-E_F.\]
Work function is related with the threshold frequency $f_0$ of photoelectric effect by
\[W=hf_0.\]
\sssc{Density or concentration}
For a type of particles, the (number) density or concentration of it is defined as the number of it per unit volume in a material.
\sssc{Number density of states of the allowed energy level per unit volume per unit energy interval}
The number density of states of the allowed energy level per unit volume per unit energy interval as a function of energy level $E$ is
\[N(E)=\frac{8\sqrt{2}\pi m_e^{3/2}}{h^3}\sqrt{E}\]
where $m_e$ is the mass of an electron.
\sssc{Occupancy probability}
The occupancy probability, the probability that a given available state will be occupied by an electron, as a function of energy level $E$, is
\[P(E)=\frac{1}{e^{\frac{E-E_F}{kT}}+1}\]
where $T$ is temperature in K and $k$ is Boltzmann's constant $8.62\times 10^{-5}$ eV/K.
\sssc{Number density of occupied states per unit volume per unit energy interval}
The number density of occupied states per unit volume per unit energy interval $N_0(E)$ is
\[N_0(E)=N(E)P(E).\]
\sssc{Finding Fermi energy}
$E_F$ can be found by solving
\[n=\int_0^{E_F}N_0(E)\dd{E}.\]
The result is
\[E_F=\qty(\frac{3}{16\sqrt{2}\pi})^{2/3}\frac{h^2}{m_e}n^{2/3}\approx\frac{0.121h^2}{m_e}n^{2/3}.\]
\sssc{Ohmic contact}
An ohmic contact is a non-rectifying electrical junction (接面), that is, a junction between two semiconductors or conductors that has a linear current–voltage (I–V) curve as with Ohm's law.
\ssc{Metals}
\sssc{Bands}
A partially filled band always exist, in which electrons can wander around and conduct current, called free electrons (or simply electrons).
\sssc{Density of free electrons in partially filled band}
Density of electrons in partially filled band $n$ is the material's mass density multiplied by the material's number of valence electrons per atom multiplied by Avogadro's number divided by the material's molar mass.
\sssc{Thermal effect}
The change of Fermi level on temperature change is small and often negligible for metals.
\sssc{Drude model of drift velocity and drift current}
In Drude model, after a typical collision, each free electron completely lose its memory of its previous velocity and moving off in a random direction. In the average time $\tau$ between collisions, the average free electron will acquire a drift speed of
\[\mathbf{v}_d=-\frac{q\tau}{m}\mb{E}\]
where $q$ is the elementary charge.

We define the free electron mobility $\mu_n$ as
\[\mu_n=\frac{q\tau}{m}\]
Suppose the free electron concentration is $n$. Then, the charge density is $-nq$ and the drift current density is
\[\mathbf{J}_d=-nq\mb{v}_d=nq\mu_n\mb{E}=\frac{nq^2\tau}{m}\mb{E}\]
Thus the resistivity $\rho$ is
\[\rho=\frac{\abs{\mathbf{E}}}{\abs{\mb{J}}}=\frac{m}{nq^2\tau}\]
which is a constant, i.e., obeying Ohm's law.
\ssc{Semiconductors}
\sssc{Bands}
For semiconductors and insulators,
\bit
\item\tb{Conduction band (導帶)}: The lowest-energy band that is not filled, in which electrons, if exist, can wander around and conduct current, called free electrons (or simply electrons). The energy level of conduction band minimum is denoted $E_C$.
\item\tb{Valence band (價帶)}: The highest-energy band that is filled, in which electrons are constrained in covalent bonds and can't wander away freely or conduct current. The energy level of valence band maximum is denoted $E_V$.
\item\tb{Band gap (能隙)}: The gap between conduction band minimum and valence band maximum. The energy difference of the gap is called band-gap energy $E_g=E_C-E_V$. For silicon, $E_g=1.12$ eV.
\eit
\sssc{Insulator}
A material with empty conduction band and high band-gap energy such that no transition of electron from valence band to conduction band occurs.
\sssc{Semiconductor}
A material whose elements of crystal structure have a valence of four in average and form a network of bonds as valence band and with low band-gap energy such that transition of electron from valence band to conduction band can occur under room temperature and that conductivity lies between that of conductors and insulators.

There are two kinds of semiconductors:
\bit
\item single-element semiconductors, such as silicon (IV) and germanium (IV), and
\item compound semiconductors, such as gallium arsenide \ce{GaAs} (III-V), aluminium gallium arsenide \ce{AlGaAs} (III-V), indium phosphide \ce{InP} (III-V), aluminium gallium indium phosphide \ce{AlInGaP} (III-V), zinc oxide \ce{ZnO2} (II-VI), and titanium dioxide \ce{TiO2} (II-VI).
\eit
Germanium was used in the fabrication of very early transistors (late 1940s, early 1950s). It was quickly supplanted by silicon, on which today’s integrated-circuit technology is almost entirely based.

Compound semiconductors are useful in special electronic circuit applications as well as in applications that involve light, such as light-emitting diodes (LEDs).
\sssc{Energy and effective mass (有效質量) of electrons}
The energy of an electron in a semiconductor is not the simple square relation of that of a free electron and wave number but a complicated periodic function. However, near an edge of a band, it can be approximated with:
\[E(k) \approx E_0 + \frac{\hbar^2 k^2}{2m^*}\]
where $E_0$ is the energy of the edge, $k$ is wave number, and $m^8$ is a constant called effective mass given by
\[m^*=\hbar^2\qty(\frac{\mathrm{d}^2E(k))}{\mathrm{d}k^2})^{-1}\]
where $\dv[2]{E(k)}{k}$ is taken at the edge. Effective mass is usually written as multiple of mass of electron.
\sssc{Effective mass of electrons at difference levels}
\bit
\item At conduction band minimum, $\dv[2]{E(k)}{k}>0$ and $m^*>0$. These electrons behave like free electrons but with mass $m^*$.
\item At valence band maximum, $\dv[2]{E(k)}{k}<0$ and $m^*<0$. These electrons move to the opposite direction of a force acting on it.
\eit
\sssc{(Electron) hole (（電子）電洞/空穴)}
A hole $h^+$ is an absence of an electron in valence band.

Let the effective mass of electrons at valence band maximum be $-m^*$ and the concentration of holes be $p$. Consider a nearly-filled valence band, the current of it equals to the current of a filled valence band, which is always zero, subtracting the current that would exist due to the electrons that would fill the holes if there are not the holes, which is in the opposite direction of the force. Thus, the current is in the direction of the force and equals to that due to particles with charge $q$ and mass $m^*$ at concentration $p$.

Thus, we define hole as a quasiparticle with charge $q$ and effective mass $m^*$, which is positive.
\sssc{Carriers (載子) of semiconductors}
Thus, there are two types of carriers in semiconductors mathematically:
\bit
\item free electron at conduction band minimum with charge $-q$ and a positive effective mass $m_n^*$, and whose properties are denoted with subscript $n$, and
\item holes at valence band maximum with charge $q$ and a positive effective mass $m_p^*$, and whose properties are denoted with subscript $p$.
\eit
The concentration of free electrons is denoted $n$; the concentration of holes is denoted $p$.
\sssc{Effective densities of states}
Available electronic states effectively participate in carrier occupation near a band edge at a given temperature, denoted as $N_C$ for conduction band minimum and $N_V$ for valence band maximum. They are governed by
\[N_C=2\qty(\frac{2\pi m_n^*kT}{h^2})^{\flatfrac{3}{2}}\]
\[N_V=2\qty(\frac{2\pi m_p^*kT}{h^2})^{\flatfrac{3}{2}}\]
\sssc{Carrier generation}
When an electron transits from valence band to conduction band, a free electron in conduction band and a hole in valence band are generated, and an energy of band-gap energy $E_g$ is consumed.

Since they are generated in pairs, an intrinsic (undoped) semiconductor has equal concentration of free electrons and holes.
\sssc{Carrier recombination (載子複合) and electroluminescence (電致發光/電激發光)}
When an electron transits from conduction band to valence band, a free electron in conduction band and a hole in valence band are cancelled, and an energy of band-gap energy $E_g$ is released in the form of a photon with frequency $f_g=\frac{E_g}{h}$, called electroluminescence.

Carrier recombination occurs naturally and at a rate proportional to the number of free electrons and holes.

Semiconductors with different band-gap energy can release light of different frequency. This effect is used in light emitting diode (LED) (發光二極體).
\sssc{Exciton (激子)}
A quasiparticle with no charge of a pair of free electron and hole in bounded state due to electrostatic force.
\sssc{Photovoltaic effect (光生伏打效應/光伏效應)/Internal photoelectric effect (內光電效應)}
Carrier generation caused by light.

Similar to the photoelectric effect, if the frequency of the incident light is greater than or equal to the material's threshold frequency $f_g=\frac{E_g}{h}$, electron-hole pairs are generated immediately; if the frequency of the incident light is less than this threshold, such pairs are never produced.

This effect is used in solar cell, photoresistor, and light-controlled relay.
\sssc{Thermal generation}
Carrier generation due to thermal energy. The higher the temperature, the more thermal generations occur.
\sssc{Thermal equilibrium}
Thermal equilibrium is the state where thermal generation rate equals to carrier recombination rate, at which
\[np=n_i^2\]
where $n_i$ denotes the number of free electrons and holes in a unit volume at a given temperature given by
\[n_i=BT^{\flatfrac{3}{2}}e^{-\frac{E_g}{2kT}}\]
where $B$ is a material-dependent parameter that is $7.3\times 10^{15}$ cm$^{-3}$K$^{-\frac{3}{2}}$ for silicon, and
\[n=N_C\exp(-\frac{E_C-E_F}{kT})\]
\[p=N_V\exp(-\frac{E_F-E_V}{kT})\]
For intrinsic semiconductor at thermal equilibrium,
\[n=p=n_i\]
\[E_F=\frac{E_C+E_V}{2}+\frac{kT}{2}\ln(\frac{N_V}{N_C})=\frac{E_C+E_V}{2}+\frac{3kT}{4}\ln(\frac{m_p^*}{m_n^*})\]
At 0 K, no thermal generation occurs, so an intrinsic semiconductor has no carrier and behaves as an insulator at thermal equilibrium. At room temperature, the carrier concentration at thermal equilibrium is still far too small for intrinsic silicon to conduct appreciable current.
\sssc{Doping (摻雜)}
Doping is the intentional introduction of impurities into an intrinsic/undoped/pure semiconductor in order to change its properties. The doped material is referred to as an extrinsic/doped semiconductor. $E_C$ and $E_V$ are approximately unchanged for light doping.
\sssc{n-typed (n 型) semiconductors or n materials}
Semiconductor with doping of elements that readily lose electrons, called donors, typically elements with valence of 5, such as phosphorus \ce{P} and arsenic \ce{As}, which introduces an increase of the concentration of free electrons, Fermi level, and conductivity. The introduced free electrons are neutralized by an equal amount of bound/bounded positive charge associated with the donor atoms even though they are often not shown in diagrams for simplicity.

If the concentration of donor atoms is $N_D$, which is usually far greater than $n_i$, the concentration of free electrons in the n-type semiconductor will be approximately
\[n_n=N_D\]
and, at thermal equilibrium,
\[p_n=\frac{n_i^2}{N_D}\]
Hence, free electrons are said to be the majority (charge) carriers, while holes are said to be the minority (charge) carriers. The minority holes are often not shown in diagrams for simplicity.
\sssc{p-typed (p 型) semiconductors or p materials}
Semiconductor with doping of elements that readily get electrons, called acceptors, typically elements with valence of 3, such as boron \ce{B} and aluminium \ce{Al}, which introduces an increase of the concentration of holes and conductivity and decrease of Fermi level. The introduced holes are neutralized by an equal amount of bound/bounded negative charge associated with the acceptor atoms even though they are often not shown in diagrams for simplicity.

If the concentration of acceptor atoms is $N_A$, which is usually far greater than $n_i$, the concentration of holes in the p-type semiconductor will be approximately
\[p_p=N_A\]
and, at thermal equilibrium,
\[n_p=\frac{n_i^2}{N_A}\]
Hence, holes are said to be the majority (charge) carriers, while free electrons are said to be the minority (charge) carriers. The minority free electrons are often not shown in diagrams for simplicity.
\ssc{Current flow in semiconductors}
\sssc{Drude model of drift velocity and drift current}
In Drude model, after a typical collision, each free electron or hole completely lose its memory of its previous velocity and moving off in a random direction. In the average time $\tau_n,\tau_p$ between collisions, the average free electron and hole will acquire a drift speed of
\[\mathbf{v}_n=-\frac{q\tau_n}{m_n^*}\mb{E}\]
\[\mathbf{v}_p=\frac{q\tau_p}{m_p^*}\mb{E}\]
We define the free electron mobility $\mu_n$ and hole mobility $\mu_p$ as
\[\mu_n=\frac{q\tau_n}{m_n^*}\]
\[\mu_p=\frac{q\tau_p}{m_p^*}\]
For silicon, $\mu_n$ is 1350 cm$^2$V$^{-1}$s$^{-1}$ and $\mu_p$ is 480 cm$^2$V$^{-1}$s$^{-1}$.

The charge densities are $-nq$ and $pq$ for free electrons and holes respectively, and the drift current densities are
\[\mathbf{J}_{S,n}=-nq\mb{v}_n=nq\mu_n\mb{E}=\frac{nq^2\tau_n}{m_n^*}\mb{E}\]
\[\mathbf{J}_{S,p}=pq\mb{v}_n=pq\mu_p\mb{E}=\frac{pq^2\tau_p}{m_p^*}\mb{E}\]
The total drift current density is
\[\mb{J}_S=\mb{J}_{S,n}+\mb{J}_{S,n}=q\qty(n\mu_n+p\mu_p)\mb{E}=q^2\qty(\frac{n\tau_n}{m_n^*}+\frac{p\tau_p}{m_p^*})\mb{E}\]
Thus the resistivity $\rho$ is
\[\rho=\frac{\abs{\mathbf{E}}}{\abs{\mb{J}_S}}=\frac{1}{q\qty(n\mu_n+p\mu_p)}=\frac{1}{q^2\qty(\frac{n\tau_n}{m_n^*}+\frac{p\tau_p}{m_p^*})}\]
which is a constant, i.e., obeying Ohm's law.
\sssc{Diffusion current}
When the concentration of a type of charge carriers differ across a piece of semiconductor, the charge carriers will diffuse from he region of high concentration to the region of low concentration, which gives rise to a net flow of charge, called diffusion current.

The concentration of charge carriers of a piece of semiconductor under a state is called concentration profile.

The diffusion current densities of free electrons and holes are
\[\mb{J}_{D,n}=-qD_n\nabla n\]
\[\mb{J}_{D,p}=qD_p\nabla p\]
where $\nabla_n,\nabla_p$ are concentration gradients, $D_n,D_p$ are constants called diffusion constants or diffusivities of free electrons and holes respectively given by:
\[D_n=\mu_nV_T\]
\[D_p=\mu_pV_T\]
called Einstein relationship, where $V_T$ is thermal voltage given by
\[V_T=\frac{kT}{q}\]
For $T=300$ K, $V_T$ is 25.9 mV. For intrinsic silicon, $D_n$ is 35 cm$^2$/s and $D_p$ is 12 cm$^2$/s.
\ssc{pn or p-n junction}
\sssc{Physical structure}
A p-n junction is a single semiconducting crystal with one end doped to form p material and the other end doped to form n material. The two types meet at a junction plane. External wire connections are made to both ends through metal contacts. If the pn junction is used as diode (二極體), these constitute the two terminals, anode, indicated by +, for p-type end and cathode, indicated by -, for n-type end, of the diode.
\sssc{(Carrier-)depletion/space-charge region/layer/zone and diffusion current}
\bit
\item Majority holes diffuse from the p region to the n region. The holes that diffuse across the junction into the n region quickly recombine with some of the majority free electrons, which results in the disappearance of some free electrons in the n region. Thus, some bound positive charge will no longer be neutralized by free electrons, and this charges are said to have been uncovered. Since recombination occurs close to the junction, there will be a region in the n region near the junction that is depleted of free electrons and contains uncovered bound positive charge.
\item Majority free electrons diffuse from the n region to the p region. The free electrons that diffuse across the junction into the p region quickly recombine with some of the majority holes, which results in the disappearance of some holes in the p region. Thus, some bound negative charge will no longer be neutralized by holes, and this charges are said to have been uncovered. Since recombination occurs close to the junction, there will be a region in the p region near the junction that is depleted of holes and contains uncovered bound negative charge.
\eit
The region on both sides of the junction, with the n side of this region positively charged and depleted of free electrons and the p side negatively charged and depleted of holes, is called depletion region.

Due to the depletion, the carrier concentrations of holes and free electrons in the whole depletion region are both negligible.

Thus, the net charge density in the region is the bound charge density.
\bit
\item In n-type part of the depletion region, the net charge density is $N_D$.
\item In p-type part of the depletion region, the net charge density is $-N_A$.
\eit
The width of the region on the two sides are not always the same. Rather, to uncover the same amount of charge, the depletion layer will extend deeper into the more lightly doped material, that is, wider into the n material if $N_A>N_D$ and wider into the p material if $N_A<N_D$, where $N_A$ and $N_D$ are the doping concentrations of the p side and n side of the junction respectively.

The charges on both sides cause an electric field $\mb{E}$ to be established in the direction from n side to p side, which results in a potential difference across the region, with the n side at a positive voltage relative to the p side, called barrier voltage $V_b$.

Both diffusions of majority carriers form currents from the n side to the p side. However, the barrier voltage acts as an energy barrier that only majority carriers with energy no less than $qV_b$ can diffuse pass the depletion region, and the others will be pushed back by the electric field.
\bit
\item For the majority holes from the p region, the pushed back holes form a drift current density $pq\mb{E}$, which is from the p side to the n side, and thus the net current density due to majority holes from the p region to the n region through the depletion region is
\[\mb{J}_{D,p}=q\qty(D_p\nabla p+p\mb{E})\]
\item For the majority free electrons from the n region, the pushed back free electrons form a drift current density $nq\mb{E}$, which is from the p side to the n side, and thus the net current density due to majority free electrons from the n region to the p region through the depletion region is
\[\mb{J}_{D,n}=q\qty(-D_n\nabla n+n\mb{E})\]
\eit
Their sum is the net current density due to majority carriers from the p region to the n region through the depletion region
\[\mb{J}_D=\mb{J}_{D,p}+\mb{J}_{D,n}=q(D_p\nabla p-D_n\nabla n+(n+p)\mb{E}),\]
The net current due to majority carriers from the p region to the n region through the depletion region is called diffusion current $I_D$, which decreases as barrier voltage increases. 
\sssc{Drift current}
\bit
\item Some of the thermally generated minority holes in the n region move toward the junction and reach the edge of the depletion region. There, they are swept across depletion region into the p side by the electric field.
\item Some of the thermally generated minority free electrons in the p region move toward the junction and reach the edge of the depletion region. There, they are swept across depletion region into the n side by the electric field.
\eit
Both currents are from the n side to the p side, and their sum is called the drift current $I_S$. Since the current is carried by thermally generated minority carriers, its value is dependent on temperature; however, it is independent of barrier voltage because the drift current is determined by the number of minority carriers that make it to the edge of the depletion region; any minority carriers that manage to get to the edge of the depletion region will be swept across.
\sssc{Current equilibrium with open-circuit terminals}
Under open-circuit condition, no external current exists; thus
\[I_D=I_S\]
This equilibrium is stable and maintained by $V_0$. If for some reason, $I_D>I_S$, more bound charges will be uncovered on both sides of the junction, which increases $V_b$ and thus decreases $I_D$, which in turns shrinks the depletion region since some of the uncovered bound charges are then covered.
\sssc{Voltage equilibrium with short-circuit terminals}
At equilibrium, the barrier voltage is balanced out by the voltages across the two metal-semiconductor junctions, whose behavior will be explained in its section. Thus, no voltage can be measured across the terminals, complying with the conservation of energy.
\sssc{(Junction) built-in voltage}
The barrier voltage $V_0$ with no external voltage applied is given by
\[V_0=V_T\ln(\frac{N_AN_T}{n_i^2}),\]
called the junction built-in voltage, which is typically in the range of 0.6 V to 0.9 V for silicon at room temperature.
\sssc{Carrier concentration under no external voltage}
Under no external voltage, which is signified by a subscript $0$,
\bit
\item at the p material edge of the depletion and further into the p material, the majority holes have concentration $p_{p0}=N_A$, and the minority free electrons have concentration $n_{p0}=\frac{n_i^2}{N_A}$;
\item at the n material edge of the depletion and further into the n material, the majority free electrons have concentration $n_{n0}=N_D$, and the minority holes have concentration $p_{n0}=\frac{n_i^2}{N_D}$.
\eit
\sssc{Charge and width of depletion region}
Let the materials be cuboid, cross sectional area of the junction be $A$, width of depletion region in $p$ side and $n$ side be $x_p$ and $x_n$ respectively, the total charge stored in depletion region in $p$ side and $n$ side be $Q_-<0$ and $Q_+>0$ respectively. Then,
\[\abs{Q_+}=qAx_nN_D=\abs{Q_-}=qAx_pN_A\]
The width of the depletion region $W=x_p+x_n$ is given by
\[W=\sqrt{\frac{2\varepsilon}{q}\qty(\frac{1}{N_A}+\frac{1}{N_D})V_b}\]
where $\varepsilon$ is the electrical permittivity of the material, which for silicon is $11.7\varepsilon_0=11.7\times 11.7\times 8.85\times 10^{-14}$ F/cm $=1.04\times 10^{-12}$ F/cm. Typical $W$ ranges from 0.1 \textmu{}m to 1 \textmu{}m.
\[x_n=W\frac{N_A}{N_A+N_D}\]
\[x_p=W\frac{N_D}{N_A+N_D}\]
The charge stored on either side of the depletion region is
\[\ba
Q_j&=\abs{Q_+}=\abs{Q_-}\\
   &=WAq\frac{N_AN_D}{N_A+N_D}\\
   &=A\sqrt{2\varepsilon q\qty(\frac{1}{N_A}+\frac{1}{N_D})V_b}
\ea\]
\sssc{On-sided junction}
A pn junction where one side is much more heavily doped than the other, that is,
\[N_A\gg N_D\qc W\approx x_n\gg x_p\]
or
\[N_D\gg N_A\qc W\approx x_p\gg x_n\]







\sssc{Junction rectifier}
A p-n junction conducts electricity better for one direction of an applied potential difference (forward biased) than for the opposite direction (back biased), and thus the device can serve as a junction rectifier.

In forward-bias connection, the p side becomes more positive and the n side becomes more negative, thus decreasing the height of the potential barrier, increasing the diffusion current, narrowing the depletion zone, and decreasing its resistivity.

In back-bias connection, the p side becomes less positive and the n side becomes less negative, thus increasing the height of the potential barrier, decreasing the diffusion current, widening the depletion zone, and increasing its resistivity.
\sssc{發光二極體（Light emitting diode, LED）}
將順向偏壓通入半導體二極體，使depletion zone極窄，使其中大量發生載子複合而電致發光。The frequency $f$ of the emitted light is
\[E_g=hf.\]

早期研發出的紅色和綠色單色發光二極體用途較受限，1992年赤崎勇、天野浩與中村修二以氮化鎵材料發明藍光 LED，始可用光的三原色調配任意所需的顏色，尤其用作照明的白光，較傳統白熾燈泡與日光燈更省電、體積更小、壽命更長，共同獲2014年諾貝爾物理獎。
\sssc{Photodiode}
Shining light on a suitably arranged p-n junction can that includes the junction. This is the basis for the photodiode.
\sssc{Junction laser}
In a forward-bias connection connected to a source, there are many electrons in the conduction band of the n-type material and many holes in the valence band of the p-type  material. Thus, there is a population inversion for the electrons; that is, there are more electrons in higher energy levels than in lower energy levels. When a single electron moves from the conduction band to the valence band, it can release its energy as a photon. This photon can stimulate a second electron to fall into the valence band, producing a second photon by stimulated emission. In this way, if the current through the junction is great enough, a chain reaction of stimulated emission events can occur and laser light can be generated. To bring this about, opposite faces of the p-n junction crystal must be flat and parallel, so that light can be reflected back and forth within the crystal. Junction lasers are built into compact disc (CD) players, where, by detecting reflections from the rotating disc, they are used to translate microscopic pits in the disc into sound. They are also much used in optical communication systems based on optical fibers.
\sssc{光伏電池（Photovoltaic cell）/太陽能電池（Solar cell）}
通常由一個或多個 p–n 接面組成，吸收光子並輸出直流電。種類如：
\begin{itemize}
	\item 單結晶矽太陽能電池：高效率但成本較高。
	\item 多晶矽太陽能電池：效率較低但成本較低。
	\item 薄膜太陽能電池：輕薄柔軟，適用於特定場景。
\end{itemize}
效率提升方法：
\begin{itemize}
	\item 減少表面反射：透明導電層和抗反射層減少光子的反射。
	\item 改進材料：新材料如有機太陽能電池和混合鈣鈦礦太陽能電池提高效率。
	\item 多接面結構：提高光吸收和轉換效率。
	\item 雷射細工：提高電池表面光陽極的粗糙度，增加光吸收。
\end{itemize}
\sssc{場效電晶體（Field-effect transistor, FET）}
FET 具有四個端子：源極（source）、閘極（gate）和汲極（drain）、體（body）/基（base）/塊體（bulk）/基板（substrate）。通過在閘極施加電壓，可以改變源極與汲極之間的電導率，從而控制電流，分為需加閘極電壓方能導通源極與汲極的增強模式（enhancement-mode），較常見；與預設導通源極與汲極，加閘極電壓可減弱或關閉它的空乏模式/減少模式（depletion-mode）。

下討論 metal-oxide-semiconductor-field-effect transistor (MOSFET)。其僅使用電子與電洞的其中一種作為電荷載子，前者稱 n-通道（n-channel），後者稱 p-通道（p-channel）。源極與汲極下方的摻雜半導體具有該種載子，在增強模式中預設不相連而斷路，但加閘極電壓可以引入載子使相連而導通；在空乏模式中預設相連而導通，但加閘極電壓可以引入相反的載子使變窄增加電阻，繼而不相連而斷路。




\ssc{Metal–semiconductor (MS or M-S) junction}
% https://en.wikipedia.org/wiki/Metal%E2%80%93semiconductor_junction


\end{document}
